% Ampliación demostrativa DF003--DF007. No modifica la fuente integral.
% Propietario común de los archivos Lean citados en estos comentarios:
% output/PAQUETE_ARTICULO_I_INTEGRACION_ESTADO_CAMPOS_20260919/.
% Residencia causal de la primera sección: después de los lectores regionales
% coinductivos y antes de emplear sus registros como entrada de la selección K.
% Residencia causal de la segunda sección: después de la construcción marcada
% Paley--Witt--Golay--Leech y antes de la reconstrucción de la estructura VOA.
% Los dos cortes expositivos conservan el mismo origen APP--TRIT--TPK.

\section[Generación regional y refinamiento]{Compatibilidad integral de la generación regional, el refinamiento cilíndrico y la incidencia}
\label{sec:df-lectores-cilindros}

La prolongación regional produce sucesiones ordenadas de bloques, junto con
sus profundidades de lectura. La proyección arquimediana y la incidencia
ternaria se calculan sobre esas mismas sucesiones. El vínculo entre ambas
realizaciones puede formularse enteramente mediante operaciones enteras:
selección por intervalos racionales, división euclídea, elevación de seis
trits, lectura de acarreo y transformación de Witt. Esta formulación permite
seguir cada coeficiente desde su bloque generador hasta su empleo posterior.

La construcción utiliza los tres lectores regionales ya obtenidos del estado
APP--TRIT--TPK. Los denotaremos por los índices
\(c\in\{\mathrm P,\mathrm E,\mathrm F\}\), correspondientes, respectivamente,
a las regiones de clausura, propagación y autoescala. Su realización real se
escribirá \(v_c\). Las identificaciones posteriores
\(v_{\mathrm P}=\pi_{\mathrm{HMT}}\),
\(v_{\mathrm E}=e_{\mathrm{HMT}}\) y
\(v_{\mathrm F}=\varphi_{\mathrm{HMT}}\) conservan ese orden causal.
La operación que emite un bloque inspecciona los intervalos racionales del
lector regional; el valor real interviene en la demostración de corrección
de dicha operación.

\subsection{Emisión posicional mediante separación racional de celdas}
\label{subsec:df-emision-racional}

% Fuente: lean/biblioteca/RegionalPublicationComposition.lean,
% produced_brackets, eventual_publication, joint_eventual_publication,
% search_terminates, stopped_cell_correct, publications_compatible.

Para cada región se dispone de dos sucesiones racionales
\(\ell_c(j),u_c(j)\), construidas por su lector, que cumplen
\[
 \ell_c(j)\leq v_c\leq u_c(j).
\]
La propiedad de separación de celdas demostrada para esos lectores afirma
que, para cada entero \(s>0\), existe \(J\) tal que, para todo \(j\geq J\),
\begin{equation}
 \lfloor s\ell_c(j)\rfloor
 =\lceil su_c(j)\rceil-1
 =\lfloor sv_c\rfloor.
 \label{eq:df-separacion-celdas}
\end{equation}
La evitación de las fronteras racionales por los valores regionales es la
hipótesis de los teoremas previos que garantiza esta separación. Aquí se
emplea su conclusión operativa, con los lectores y los intervalos que la
producen.

\begin{definition}[Índice de separación y prefijo emitido]
Para \(s>0\), sea
\[
 j_c(s)=\min\{j\in\mathbb N:
 \lfloor s\ell_c(j)\rfloor=\lceil su_c(j)\rceil-1\}.
\]
Para una base entera \(b\geq2\) y una profundidad \(n\geq0\), se define
\begin{equation}
 P_c(b,n)=\lfloor b^n\ell_c(j_c(b^n))\rfloor.
 \label{eq:df-publicacion-prefijo}
\end{equation}
El prefijo contiene la parte entera y las primeras \(n\) posiciones en base
\(b\); por consiguiente, \(P_c(b,0)\) es la parte entera regional.
\end{definition}

\begin{theorem}[Terminación y corrección de la emisión]
\label{thm:df-emision-correcta}
El índice \(j_c(s)\) existe para todo \(s>0\). Los prefijos emitidos cumplen
\begin{equation}
 P_c(b,n)=\lfloor b^n v_c\rfloor,
 \qquad
 P_c(b,n)\leq b^n v_c<P_c(b,n)+1.
 \label{eq:df-prefijo-correcto}
\end{equation}
Para cada escala \(s\) puede elegirse además un único índice de suficiencia
\(J(s)\) válido simultáneamente para las tres regiones.
\end{theorem}

\begin{proof}
La igualdad eventual de \eqref{eq:df-separacion-celdas} hace no vacío el
conjunto que define \(j_c(s)\). Su mínimo existe por buen orden de
\(\mathbb N\). En cualquier índice donde se produzca la igualdad, la cota
inferior racional proporciona
\(\lfloor s\ell_c(j)\rfloor\leq\lfloor sv_c\rfloor\).
La cota superior, junto con la evitación de una frontera de celda, da
\(\lfloor sv_c\rfloor\leq\lceil su_c(j)\rceil-1\).
Al coincidir ambos extremos enteros, el entero intermedio coincide con
ellos. Esto demuestra \eqref{eq:df-prefijo-correcto}. Para la afirmación
conjunta basta tomar el máximo de los tres índices de suficiencia obtenidos
para la misma escala. El índice conjunto controla simultáneamente las tres
emisiones, mientras que cada región conserva sus propios intervalos.
\end{proof}

\begin{proposition}[Compatibilidad exacta de los prefijos]
\label{prop:df-prefijos-compatibles}
Para \(n,m\geq0\),
\begin{equation}
 \left\lfloor\frac{P_c(b,n+m)}{b^m}\right\rfloor=P_c(b,n).
 \label{eq:df-truncamiento}
\end{equation}
En particular, el bloque siguiente
\(d_c(b,n)=P_c(b,n+1)\bmod b\) satisface
\begin{equation}
 P_c(b,n+1)=bP_c(b,n)+d_c(b,n),\qquad 0\leq d_c(b,n)<b.
 \label{eq:df-paso-posicional}
\end{equation}
\end{proposition}

\begin{proof}
Escribamos \(b^n v_c=N+\theta\), con
\(N=P_c(b,n)\) y \(0\leq\theta<1\). Entonces
\(P_c(b,n+m)=b^mN+\lfloor b^m\theta\rfloor\), y el segundo sumando
pertenece a \(\{0,\ldots,b^m-1\}\). La división euclídea por \(b^m\)
recupera \(N\). La especialización \(m=1\) produce
\eqref{eq:df-paso-posicional} y su cota residual.
\end{proof}

\subsection{Bloques regionales de seis trits a profundidad arbitraria}
\label{subsec:df-bloques-arbitrarios}

% Fuente: deltas/integracion/JointRegionalFrontier.lean,
% publish_base729_eq_base3, generated_block_from_longer_prefix,
% generated_block_eq_finite, generated_signature_eq_finite.

La igualdad entera \(729=3^6\) determina la agrupación de seis posiciones
ternarias en un bloque. Para \(t\geq0\), definimos
\begin{equation}
 u_c(t)=P_c(729,t+1)\bmod729,
 \qquad
 B_{c,j}(t)=
 \left\lfloor\frac{u_c(t)}{3^{5-j}}\right\rfloor\bmod3,
 \quad 0\leq j<6.
 \label{eq:df-bloque-y-banda}
\end{equation}
De este modo, \(B_c(t)\) es una palabra ordenada de seis trits y
\begin{equation}
 u_c(t)=\sum_{j=0}^{5}B_{c,j}(t)3^{5-j}.
 \label{eq:df-elevacion-banda}
\end{equation}
La palabra y su elevación entera son dos representaciones mutuamente
recuperables del mismo bloque emitido.

\begin{theorem}[Estabilidad de extracción a cualquier profundidad]
\label{thm:df-bloque-estable}
Se cumplen, para todo \(n\),
\[
 P_c(729,n)=P_c(3,6n).
\]
Si \(t<n\), entonces
\begin{equation}
 u_c(t)=
 \left\lfloor\frac{P_c(729,n)}{729^{n-t-1}}\right\rfloor\bmod729.
 \label{eq:df-bloque-prefijo-largo}
\end{equation}
En consecuencia, cualquier ampliación del prefijo preserva todos los
bloques previamente emitidos y todas las firmas que se calculan sobre ellos.
\end{theorem}

\begin{proof}
La primera igualdad resulta de que ambos prefijos se separan a la misma
escala, \(729^n=3^{6n}\), y del
teorema~\ref{thm:df-emision-correcta}. Para la segunda aplicamos
\eqref{eq:df-truncamiento} a las profundidades \(n\) y \(t+1\), y después
tomamos el residuo módulo \(729\). La extracción ternaria de
\eqref{eq:df-bloque-y-banda} es determinista. Cualquier función explícita
de esas dieciocho coordenadas conserva, por tanto, el mismo resultado al
ampliar el prefijo.
\end{proof}

La construcción está cuantificada sobre \(t\in\mathbb N\). Una tabla que
materialice cien bloques constituye una evaluación finita de esta
construcción; la fórmula \eqref{eq:df-bloque-prefijo-largo} determina
exactamente su compatibilidad con cualquier extensión posterior.

\begin{example}[Primer bloque conjunto]
La evaluación de los lectores regionales seleccionados da las palabras
\[
 B_{\mathrm P}(0)=010211,\qquad
 B_{\mathrm E}(0)=201101,\qquad
 B_{\mathrm F}(0)=121200.
\]
Su elevación por \eqref{eq:df-elevacion-banda} produce
\[
 u_{\mathrm P}(0)=103,\qquad
 u_{\mathrm E}(0)=523,\qquad
 u_{\mathrm F}(0)=450.
\]
Por ejemplo,
\(201101_3=2\cdot243+27+9+1=523\).
Las seis posiciones, incluidos los ceros iniciales, se conservan en la
representación de banda. El cero inicial de \(010211\) tiene una posición
definida aunque la escritura decimal de \(103\) carezca de él.
\end{example}

\subsection{Registros de transición y reconstrucción de las elevaciones lineales}
\label{subsec:df-registros-transicion}

% Fuentes: deltas/integracion/RegionalTransitionRecords.lean y
% GeneratedTransitionRecords.lean; lean/biblioteca/TPKLifts.lean.
% Los enlaces reconstruidos son R(0)->R(1) y R(2)->R(3).

Sobre \(\mathbb F_3\), reunimos dos emisiones consecutivas de cada región
en la matriz
\begin{equation}
 R(t)=
 \begin{pmatrix}
 B_{\mathrm P}(t)\\ B_{\mathrm P}(t+1)\\
 B_{\mathrm E}(t)\\ B_{\mathrm E}(t+1)\\
 B_{\mathrm F}(t)\\ B_{\mathrm F}(t+1)
 \end{pmatrix}.
 \label{eq:df-registro-matricial}
\end{equation}
Esta definición fija tanto el orden regional como el orden temporal. Los
cuatro primeros registros son
\[
 X_0=R(0),\quad Y_0=R(1),\quad X_1=R(2),\quad Y_1=R(3).
\]
Sus valores, obtenidos por extracción de los prefijos regionales, son
\begin{align*}
 X_0&=\begin{pmatrix}
0&1&0&2&1&1\\0&1&2&2&2&2\\2&0&1&1&0&1\\
1&2&1&2&2&1\\1&2&1&2&0&0\\1&1&2&2&0&2
\end{pmatrix},&
 Y_0&=\begin{pmatrix}
0&1&2&2&2&2\\0&1&0&2&1&1\\1&2&1&2&2&1\\
1&0&2&0&1&1\\1&1&2&2&0&2\\1&2&1&0&2&0
\end{pmatrix},\\[1ex]
 X_1&=\begin{pmatrix}
0&1&0&2&1&1\\0&0&2&1&1&1\\1&0&2&0&1&1\\
0&1&2&2&2&2\\1&2&1&0&2&0\\0&1&0&2&1&0
\end{pmatrix},&
 Y_1&=\begin{pmatrix}
0&0&2&1&1&1\\1&1&0&2&2&1\\0&1&2&2&2&2\\
1&0&2&0&1&1\\0&1&0&2&1&0\\0&1&0&2&0&0
\end{pmatrix}.
\end{align*}

\begin{theorem}[Reconstrucción única de los dos enlaces iniciales]
\label{thm:df-elevaciones-reconstruidas}
Las ecuaciones \(X_0L_0=Y_0\) y \(X_1L_1=Y_1\) tienen, cada una, una
única solución en \(M_6(\mathbb F_3)\). Estas soluciones son
\begin{align*}
 L_0&=\begin{pmatrix}
2&2&2&1&2&1\\2&1&2&2&1&1\\1&0&0&1&0&2\\
0&0&1&1&1&0\\2&2&2&0&2&1\\2&1&2&1&0&0
\end{pmatrix},&
 L_1&=\begin{pmatrix}
2&1&2&0&2&2\\1&2&1&1&2&0\\1&2&1&1&1&2\\
0&1&0&0&2&1\\2&0&2&1&0&1\\0&2&2&2&1&1
\end{pmatrix}.
\end{align*}
Ambas matrices son invertibles.
\end{theorem}

\begin{proof}
Las matrices
\begin{align*}
 A_0&=\begin{pmatrix}
1&1&2&0&1&2\\0&0&1&0&0&1\\2&1&2&1&0&1\\
0&2&0&1&0&2\\2&1&1&0&2&2\\2&1&1&1&1&2
\end{pmatrix},&
 A_1&=\begin{pmatrix}
1&0&1&2&0&0\\0&0&1&1&2&2\\0&1&2&0&1&1\\
1&1&0&2&0&0\\1&1&2&1&1&2\\1&0&0&0&0&2
\end{pmatrix}
\end{align*}
satisfacen \(A_iX_i=X_iA_i=I\), por multiplicación en \(\mathbb F_3\).
Por tanto, las soluciones están forzadas por \(L_i=A_iY_i\); la
multiplicación produce las dos matrices mostradas. Sus inversas son
\begin{align*}
 L_0^{-1}&=\begin{pmatrix}
1&2&0&2&0&2\\2&0&2&2&0&0\\2&1&2&0&2&0\\
1&0&0&0&2&0\\0&2&1&1&2&0\\2&2&2&2&2&2
\end{pmatrix},&
 L_1^{-1}&=\begin{pmatrix}
2&0&2&1&2&1\\2&1&0&0&2&0\\2&0&2&2&0&2\\
2&0&0&1&1&0\\1&1&1&1&1&0\\2&0&1&2&2&0
\end{pmatrix}.
\end{align*}
La verificación bilateral \(L_iL_i^{-1}=L_i^{-1}L_i=I\) demuestra la
recuperación exacta de cada palabra transformada. Todos los productos
utilizan residuos \(0,1,2\), con su interpretación en \(\mathbb F_3\).
\end{proof}

La lectura causal de este resultado es precisa: las emisiones regionales
producen \(R(t)\); los registros \(R(0),\ldots,R(3)\) determinan las dos
elevaciones anteriores; sus inversas recuperan las palabras transformadas.
La continuidad para todo \(t\) pertenece a los lectores coinductivos y al
teorema~\ref{thm:df-bloque-estable}. La reconstrucción matricial aquí
demostrada tiene por dominio los dos enlaces especificados en su enunciado.

\subsection{Refinamiento simultáneo en las bases ternaria y decimal}
\label{subsec:df-cilindros-mixtos}

% Fuente: deltas/integracion/RegionalCylinderTransition.lean.

Una emisión ternaria de seis trits y una emisión decimal de tres cifras
constituyen dos lecturas de distinta base. Para conservar su relación se
registra el estado
\[
 s=(n,k,N,D),
\]
donde \(N\) es un prefijo a profundidad \(n\) en base \(729\), y \(D\) es
un prefijo a profundidad \(k\) en base \(1000\). Sus cilindros
arquimedianos son
\begin{equation}
 I_{729}(s)=\left[\frac{N}{729^n},\frac{N+1}{729^n}\right),\qquad
 I_{1000}(s)=\left[\frac{D}{1000^k},\frac{D+1}{1000^k}\right).
 \label{eq:df-cilindros}
\end{equation}
El extremo superior semiabierto asigna cada valor a una única celda de una
base y una profundidad dadas.

\begin{definition}[Defectos enteros de compatibilidad]
Definimos
\begin{align}
 A(s)&=N1000^k-D729^n,\label{eq:df-defecto-inferior}\\
 B(s)&=(N+1)1000^k-D729^n=A(s)+1000^k.
 \label{eq:df-defecto-superior}
\end{align}
El estado es compatible cuando
\begin{equation}
 A(s)<729^n,\qquad B(s)>0.
 \label{eq:df-compatibilidad}
\end{equation}
\end{definition}

\begin{lemma}[Criterio entero de intersección]
Los cilindros de \eqref{eq:df-cilindros} tienen intersección no vacía si y
sólo si se cumple \eqref{eq:df-compatibilidad}.
\end{lemma}

\begin{proof}
Dos intervalos semiabiertos \([a,b)\) y \([c,d)\), con \(a<b\) y
\(c<d\), se cortan exactamente cuando \(a<d\) y \(c<b\). Aplicando esta
condición a \eqref{eq:df-cilindros} y multiplicando por los denominadores
positivos, obtenemos
\(N1000^k<(D+1)729^n\) y
\(D729^n<(N+1)1000^k\). Son las dos desigualdades de
\eqref{eq:df-compatibilidad}.
\end{proof}

\begin{definition}[Actualización de profundidades y prefijos]
Para \(0\leq u<729\) y \(0\leq d<1000\), definimos
\begin{align*}
 T_u(n,k,N,D)&=(n+1,k,729N+u,D),\\
 T_{u,d}(n,k,N,D)&=(n+1,k+1,729N+u,1000D+d).
\end{align*}
La primera operación refina únicamente el cilindro ternario. La segunda
refina ambas lecturas y conserva las dos profundidades en el estado.
\end{definition}

\begin{proposition}[Recurrencias exactas de frontera]
\label{prop:df-rec-frontera}
Las actualizaciones anteriores satisfacen
\begin{align}
 A(T_u s)&=729A(s)+u1000^k,\label{eq:df-rec-ternaria}\\
 A(T_{u,d}s)&=729000A(s)+u1000^{k+1}-d729^{n+1}.
 \label{eq:df-rec-conjunta}
\end{align}
\end{proposition}

\begin{proof}
Para la primera identidad,
\[
 (729N+u)1000^k-D729^{n+1}
 =729(N1000^k-D729^n)+u1000^k.
\]
Para la segunda,
\begin{align*}
 &(729N+u)1000^{k+1}-(1000D+d)729^{n+1}\\
 &\qquad=729000(N1000^k-D729^n)
      +u1000^{k+1}-d729^{n+1}.
\end{align*}
Ambas son identidades enteras anteriores a cualquier aproximación real.
\end{proof}

\begin{theorem}[Refinamiento de los cilindros regionales generados]
\label{thm:df-refinamiento-generado}
Sea
\[
 s_c(n,k)=(n,k,P_c(729,n),P_c(1000,k)).
\]
Entonces \(s_c(n,k)\) es compatible para cualesquiera \(n,k\). Además,
\begin{align*}
 T_{u_c(n)}s_c(n,k)&=s_c(n+1,k),\\
 T_{u_c(n),d_c(1000,k)}s_c(n,k)&=s_c(n+1,k+1).
\end{align*}
La división euclídea de los prefijos refinados recupera exactamente los
prefijos anteriores. Para todo \(a,b\geq0\),
\[
 \frac{P_c(729,n+a)-P_c(729,n+a)\bmod729^a}{729^a}
 =P_c(729,n),
\]
y se cumple la identidad análoga en base \(1000\).
\end{theorem}

\begin{proof}
El teorema~\ref{thm:df-emision-correcta} sitúa \(v_c\) simultáneamente en
los dos intervalos \eqref{eq:df-cilindros}; el criterio de intersección da
la compatibilidad. Las identidades de actualización son
\eqref{eq:df-paso-posicional} en cada base. La recuperación se obtiene de
\eqref{eq:df-truncamiento}. Por tanto, la recurrencia del defecto y la
recuperación del prefijo son consecuencias de una misma actualización.
\end{proof}

\begin{example}[Primer par de cilindros regionales]
Para \(n=k=1\), la evaluación de las tres regiones da
\[
\begin{array}{c|rr|rr}
 c&N&D&A&B\\\hline
 \mathrm P&2290&3141&211&1211\\
 \mathrm E&1981&2718&-422&578\\
 \mathrm F&1179&1618&-522&478
\end{array}
\]
En las tres filas \(A<729\) y \(B>0\). La primera columna de prefijos
se descompone como \(2290=3\cdot729+103\),
\(1981=2\cdot729+523\) y \(1179=729+450\): reaparecen exactamente los
bloques de seis trits ya emitidos. Los signos diferentes de \(A\) expresan
la posición relativa de los dos extremos inferiores, sin alterar la
pertenencia común del valor regional.
\end{example}

\subsection{Selección del bloque siguiente mediante el lector regional}

La compatibilidad de dos intervalos expresa una relación entre lecturas.
La determinación del siguiente bloque utiliza además el intervalo racional
producido por la región. Esta operación puede formularse sin eliminar la
historia de profundidad.

\begin{definition}[Admisibilidad de la prolongación regional]
Fijados \(c,n\), sea \(s=729^{n+1}\) y \(j=j_c(s)\). Un residuo
\(u\in\{0,\ldots,728\}\) es admisible cuando
\begin{equation}
 \lfloor s\ell_c(j)\rfloor
 \leq729P_c(729,n)+u
 \leq\lceil su_c(j)\rceil-1.
 \label{eq:df-admisibilidad}
\end{equation}
\end{definition}

\begin{theorem}[Unicidad de la prolongación compatible con el lector]
\label{thm:df-prolongacion-unica}
La condición \eqref{eq:df-admisibilidad} equivale a \(u=u_c(n)\).
Por tanto, existe exactamente una prolongación admisible de cada región
a cada profundidad.
\end{theorem}

\begin{proof}
En el índice de separación, los dos extremos enteros de
\eqref{eq:df-admisibilidad} coinciden con \(P_c(729,n+1)\).
La condición equivale así a
\(729P_c(729,n)+u=P_c(729,n+1)\).
Por \eqref{eq:df-paso-posicional}, el lado derecho es
\(729P_c(729,n)+u_c(n)\). La cancelación del término común da la
igualdad de residuos. Recíprocamente, esa igualdad satisface la condición.
\end{proof}

La dependencia efectiva de esta selección queda completamente determinada:
región seleccionada, lector racional, profundidad, índice de separación,
prefijo anterior y residuo siguiente. El defecto \(A\) resume una relación
de frontera; el estado \((n,k,N,D)\) y los intervalos regionales conservan
los datos empleados por la actualización.

\subsection{Firma ternaria, acarreo, incidencia de Witt y reloj de lectura}
\label{subsec:df-firma-conjunta}

% Fuentes: lean/regional/N69RegionalSignature.lean,
% deltas/integracion/JointRegionalFrontier.lean y JointCylinderSignature.lean.

Fijada una profundidad \(t\), escribamos \(B_{c,j}=B_{c,j}(t)\). Para
cada columna \(j\) se calculan los siguientes enteros y residuos:
\begin{align}
 S_j&=B_{\mathrm P,j}+B_{\mathrm E,j}+B_{\mathrm F,j},
 &q_j&=S_j\bmod3,
 &c_j&=\left\lfloor S_j/3\right\rfloor,
 \label{eq:df-firma-qc}\\
 a_j&=(B_{\mathrm P,j}+B_{\mathrm E,j}-B_{\mathrm F,j})\bmod3,
 &w_j&=\sum_c\mathbf1_{B_{c,j}\ne0},
 &\bar w_j&=w_j\bmod3.
 \label{eq:df-firma-axial}
\end{align}
El residuo de la expresión axial se toma en los enteros. Una implementación
con naturales lo calcula como
\((B_{\mathrm P,j}+B_{\mathrm E,j}+3-B_{\mathrm F,j})\bmod3\), que conserva
exactamente ese residuo.

La matriz de incidencia ternaria utilizada por el lector es
\begin{equation}
 W=\begin{pmatrix}
0&1&1&1&1&1\\
1&0&1&1&2&2\\
1&1&0&2&1&2\\
1&1&2&0&2&1\\
1&2&1&2&0&1\\
1&2&2&1&1&0
\end{pmatrix}.
\label{eq:df-matriz-witt}
\end{equation}
Para cada región se conservan dos cargas residuales:
\begin{equation}
 v_c^{\mathrm{res}}=\sum_j B_{c,j}\bmod3,
 \qquad
 v_c^{\mathrm W}=\sum_j(B_cW)_j\bmod3.
 \label{eq:df-cargas-residuales}
\end{equation}
La firma completa del lector es
\(\mathcal S(B)=(q,a,c,w,\bar w,v^{\mathrm{res}},v^{\mathrm W})\).
Su estructura conserva por separado el residuo de columna, el acarreo, el
contraste axial, el peso de soporte y las dos lecturas de carga.

\begin{theorem}[Recuperación exacta y cotas de la firma]
\label{thm:df-recuperacion-firma}
Para cada columna y cada región se cumplen
\begin{align}
 S_j&=q_j+3c_j,&0\leq q_j,a_j&<3,&0\leq c_j&\leq2,
 \label{eq:df-recuperacion-total}\\
 B_{\mathrm F,j}&=2(q_j+3-a_j)\bmod3,&0\leq w_j&\leq3,
 &0\leq\bar w_j&<3.
 \label{eq:df-recuperacion-eje}
\end{align}
Además,
\begin{equation}
 v_c^{\mathrm W}
 =\left(\sum_j\bigl((B_cW)_j\bmod3\bigr)\right)\bmod3.
 \label{eq:df-witt-reduccion}
\end{equation}
Las recuperaciones enunciadas determinan el total de cada columna y su
coordenada de autoescala; el registro regional conserva adicionalmente
las tres bandas ordenadas.
\end{theorem}

\begin{proof}
La identidad \eqref{eq:df-recuperacion-total} es la división euclídea de
\(S_j\) por \(3\). Cada sumando de \(S_j\) está entre \(0\) y \(2\),
de modo que \(0\leq S_j\leq6\) y \(c_j\leq2\). En \(\mathbb F_3\),
la resta de las dos congruencias que definen \(q_j\) y \(a_j\) da
\(q_j-a_j=2B_{\mathrm F,j}\). Como el inverso de \(2\) es \(2\), se
obtiene \eqref{eq:df-recuperacion-eje}; el sumando \(3\) permite
escribirla con naturales. Las cotas de \(w_j\) resultan de sumar tres
indicadores. Finalmente, reducir una suma entera módulo \(3\) equivale
a reducir cada sumando y volver a reducir la suma, lo que prueba
\eqref{eq:df-witt-reduccion}.
\end{proof}

\begin{example}[Firma del primer bloque conjunto]
Las tres bandas iniciales producen
\begin{align*}
 q&=(0,0,2,2,1,2),&a&=(1,2,0,1,1,2),\\
 c&=(1,1,0,1,0,0),&w&=(2,2,2,3,1,2),\\
 v^{\mathrm{res}}&=(2,2,0),&v^{\mathrm W}&=(2,1,1).
\end{align*}
La cuarta columna tiene total \(2+1+2=5\); su firma registra
\(q_3=2\) y \(c_3=1\), por lo que \(q_3+3c_3=5\).
La recuperación axial da
\(2(2+3-1)\bmod3=2\), precisamente la coordenada de autoescala.
Las imágenes de Witt, reducidas componente a componente, son
\[
 (2,0,2,1,1,2),\qquad
 (0,0,0,2,0,2),\qquad
 (2,1,1,2,1,0).
\]
Su suma residual por fila produce las tres cargas \((2,1,1)\).
\end{example}

\begin{definition}[Reloj entero de lectura decimal]
El índice de transición \(t\) y la profundidad decimal de publicación se
relacionan mediante
\begin{equation}
 \kappa(t)=\max\{k\in\mathbb N:1000^k\leq3^{6(t+1)}\}.
 \label{eq:df-reloj}
\end{equation}
\end{definition}

\begin{proposition}[Caracterización y monotonía del reloj]
\label{prop:df-reloj}
El entero \(\kappa(t)\) queda determinado unívocamente por
\begin{equation}
 1000^{\kappa(t)}\leq729^{t+1}<1000^{\kappa(t)+1}.
 \label{eq:df-reloj-cotas}
\end{equation}
La función \(\kappa\) es monótona y satisface
\(\kappa(20)=\kappa(21)=20\).
\end{proposition}

\begin{proof}
Las potencias de \(1000\) son estrictamente crecientes y no acotadas;
existe, por tanto, un único intervalo consecutivo de potencias que
contiene \(729^{t+1}\). La monotonía de esta última expresión demuestra
la de \(\kappa\). Las dos evaluaciones finales son las comparaciones
enteras
\[
 1000^{20}\leq729^{21}<1000^{21},\qquad
 1000^{20}\leq729^{22}<1000^{21}.
\]
Estas comparaciones se efectúan con potencias enteras exactas.
\end{proof}

La igualdad de profundidad decimal en dos transiciones distintas explica
la necesidad de conservar ambos índices: el reloj de lectura puede
permanecer constante mientras avanza la banda ternaria y se actualiza su
firma. La información cronológica se encuentra en el registro ordenado
de las transiciones.

\subsection{Teorema de composición del bloque, el cilindro y la firma}

\begin{theorem}[Compatibilidad conjunta de las lecturas regionales]
\label{thm:df-composicion-conjunta}
Para cada profundidad \(n\), existe una única terna de bloques
\((u_{\mathrm P},u_{\mathrm E},u_{\mathrm F})\) que satisface la
admisibilidad regional \eqref{eq:df-admisibilidad} en las tres regiones.
Esta terna es \((u_{\mathrm P}(n),u_{\mathrm E}(n),u_{\mathrm F}(n))\).
Simultáneamente:
\begin{enumerate}
\item cada bloque actualiza su cilindro por las recurrencias
\eqref{eq:df-rec-ternaria} y \eqref{eq:df-rec-conjunta};
\item sus seis trits producen la firma
\(\mathcal S(B_{\mathrm P}(n),B_{\mathrm E}(n),B_{\mathrm F}(n))\);
\item los totales de columna se recuperan como \(q+3c\), y las cargas
de Witt se obtienen de las mismas bandas mediante
\eqref{eq:df-witt-reduccion};
\item todo prefijo de mayor profundidad devuelve exactamente estos
bloques, cilindros truncados y firmas.
\end{enumerate}
\end{theorem}

\begin{proof}
Aplicamos el teorema~\ref{thm:df-prolongacion-unica} a cada uno de los tres
índices regionales. La unicidad coordenada a coordenada da la unicidad
conjunta. Sustituimos esos mismos residuos en las operaciones
\(T_u,T_{u,d}\), lo que permite aplicar el
teorema~\ref{thm:df-refinamiento-generado}. La aplicación determinista de
\eqref{eq:df-bloque-y-banda} produce las bandas utilizadas por
\eqref{eq:df-firma-qc}--\eqref{eq:df-cargas-residuales}. Las identidades
de recuperación son las del teorema~\ref{thm:df-recuperacion-firma}.
Finalmente, el teorema~\ref{thm:df-bloque-estable} y el truncamiento de
cada base demuestran la estabilidad a mayor profundidad. En cada paso se
ha conservado la identidad del bloque emitido; ninguna de las dos
lecturas requiere seleccionar otro bloque.
\end{proof}

La conclusión es una compatibilidad operacional entre la contracción
cilíndrica y la conservación de incidencia. Ambas utilizan el mismo
registro generado, mantienen las profundidades y admiten recuperación
entera. Este resultado es la forma explícita del enlace
\[
 \text{bloque regional emitido}
 \longrightarrow
 \begin{cases}
 \text{cilindro y frontera entera},\\
 \text{banda, firma, acarreo y carga de Witt}.
 \end{cases}
\]
El desarrollo posterior de la incidencia y de los campos emplea la
segunda lectura, conservando la procedencia común con la primera.

\section[Campos sobre el retículo de incidencia]{Generación y localidad de los campos sobre el retículo de incidencia}
\label{sec:df-campos-reticulares}

La construcción reticular se realiza después de la selección del registro
dodecafásico y de sus incidencias. El origen marcado, el producto entero y
la forma de paridad que utiliza el espacio de estados proceden de esa
construcción. A continuación se explicitan las operaciones que producen
los campos, sus coeficientes y sus relaciones de localidad. El mismo
retículo interviene en la generación de estados, en los productos sobre
el vacío y en los conmutadores de todos los modos enteros.

\subsection{Residencia del origen marcado y persistencia de la incidencia}

% Fuentes: lean/incidencia/SelectedRegionalIncidence.lean;
% deltas/union/ConcreteNativeFourPlusOne.lean y SelectedExceptionalChain.lean.
% El antecedente SelectedFromRegionalInputs conserva en su propietario la
% interfaz S8 no ordenada y Lean.ofReduceBool del selector terminal. Este
% capítulo desarrolla sus salidas; no declara eliminadas esas dependencias.

Escribamos \(K_i\), \(0\leq i<12\), para las coordenadas del registro
seleccionado. Las dos lecturas de soporte son
\begin{equation}
 T_K=\{i:K_i\geq729\},\qquad
 H_K=\{i:p_i+e_i-f_i-K_i<0\},
 \label{eq:df-soportes-K}
\end{equation}
donde \(p_i,e_i,f_i\) son los bloques regionales utilizados por el lector
de preacarreo. El soporte superior utiliza el umbral entero \(729\);
el soporte negativo utiliza el signo del preacarreo antes de reducirlo.
Se conservan así dos operaciones diferentes sobre el mismo registro.

La selección ya establecida da
\[
 K=(234,543,140,729,659,824,621,58,914,794,146,601),
\]
como evaluación del registro, y produce la bandera
\(T_K\subset H_K\), con \(|T_K|=4\), \(|H_K|=6\).
Las incidencias regionales de propagación y autoescala, junto con el
soporte APP, determinan el conjunto
\[
 O_K=(E_{\mathrm{inc}}\cap F_{\mathrm{inc}})
       \setminus(H_K\cup A_{\mathrm{inc}}).
\]
La construcción de incidencia demuestra que \(O_K\) es un singleton.
Su único elemento, \(o\), fija el origen marcado que se conserva en
todo lo que sigue.

En la realización concreta de Golay, el soporte \(T_K\) corresponde a
\(\{4,6,8,14\}\). Las cuatro hexadas incidentes a ese soporte tienen
cuatro elevaciones octádicas; la estrella octádica completa tiene cinco
elementos. Su diferencia es una única octada adicional. El soporte
\(H_K\), por su parte, tiene una elevación octádica única en la misma
realización. Estas propiedades proceden de la realización residual
concreta y mantienen el origen regional del soporte seleccionado.

El vecino reticular marcado \(\Lambda_o\) obtenido de esa incidencia
tiene rango \(24\), forma bilineal entera y par, autodualidad integral
y mínimo cuadrático \(4\). El vector radial utilizado por la construcción
marcada tiene cuadrado \(54\). Denotaremos por
\(\langle x,y\rangle\in\mathbb Z\) la forma entera de \(\Lambda_o\).
Los teoremas de esta sección utilizan precisamente ese retículo y esa
forma ya construidos.

\subsection{Cociclo de paridad y álgebra reticular torcida}
\label{subsec:df-cociclo}

% Fuentes: deltas/voa/TriangularSignCocycle.lean,
% WittLatticeCocycle.lean y WittComplexAlgebra.lean.

Elijamos una base integral \((b_i)_{i=1}^{r}\) de \(\Lambda_o\), con
\(r=24\). Esta elección presenta las coordenadas de un retículo ya
obtenido. Su matriz de Gram es
\(G_{ij}=\langle b_i,b_j\rangle\). Para
\(x=\sum_i x_i b_i\), \(y=\sum_i y_i b_i\), definimos
\begin{equation}
 \tau(x,y)=\sum_{i<j}x_iG_{ij}y_j\pmod2,
 \qquad \varepsilon(x,y)=(-1)^{\tau(x,y)}.
 \label{eq:df-cociclo}
\end{equation}

\begin{proposition}[Identidad cocíclica y conmutador de paridad]
\label{prop:df-cociclo}
Se cumplen
\begin{align}
 \varepsilon(x,y)\varepsilon(x+y,z)
 &=\varepsilon(y,z)\varepsilon(x,y+z),\label{eq:df-identidad-cociclica}\\
 \varepsilon(x,y)\varepsilon(y,x)
 &=(-1)^{\langle x,y\rangle},\label{eq:df-conmutador-cociclo}\\
 \varepsilon(0,x)&=\varepsilon(x,0)=1.
\end{align}
\end{proposition}

\begin{proof}
La bilinealidad de \(\tau\) en \(\mathbb F_2\) da
\(\tau(x,y)+\tau(x+y,z)=\tau(y,z)+\tau(x,y+z)\), y al aplicar
\(a\mapsto(-1)^a\) se obtiene \eqref{eq:df-identidad-cociclica}.
La simetría de \(G\), junto con el carácter par de su diagonal, permite
reunir las dos sumas triangulares:
\[
 \tau(x,y)+\tau(y,x)
 =\sum_{i,j}x_iG_{ij}y_j
 =\langle x,y\rangle\pmod2.
\]
Esto prueba \eqref{eq:df-conmutador-cociclo}. Las identidades con cero
resultan directamente de \eqref{eq:df-cociclo}.
\end{proof}

El álgebra compleja torcida \(\mathbb C_{\varepsilon}[\Lambda_o]\)
consiste en las combinaciones lineales finitas de símbolos
\(\mathbf e^x\), con producto
\begin{equation}
 \mathbf e^x\mathbf e^y=\varepsilon(x,y)\mathbf e^{x+y}.
 \label{eq:df-producto-torcido}
\end{equation}
La identidad cocíclica demuestra la asociatividad y
\(\mathbf e^0\) es la unidad. La escritura \(\mathbf e^x\) designa el
elemento básico de carga \(x\), y conserva la etiqueta reticular completa.

\subsection{Espacio oscilatorio y representación de Heisenberg}
\label{subsec:df-osciladores}

% Fuente: deltas/voa/LatticeOscillatorFock.lean.

Sea \(\mathfrak h=\mathbb C\otimes_{\mathbb Z}\Lambda_o\). El espacio
oscilatorio algebraico es
\[
 \mathfrak h_-=
 \bigoplus_{n\geq1}\mathfrak h\,t^{-n},
 \qquad
 M(1)=\operatorname{Sym}_{\mathbb C}(\mathfrak h_-).
\]
El soporte de cada vector es finito, mientras que el conjunto de frecuencias
disponibles comprende todos los enteros positivos. El espacio de estados
que conserva simultáneamente ocupación oscilatoria y carga reticular es
\begin{equation}
 V_{\Lambda_o}=M(1)\otimes_{\mathbb C}
                  \mathbb C_{\varepsilon}[\Lambda_o],
 \qquad \mathbf1=1\otimes\mathbf e^0.
 \label{eq:df-espacio-estados}
\end{equation}

Para \(x\in\Lambda_o\) y \(n\geq1\), el operador \(x(-n)\) multiplica
por el generador \(x\,t^{-n}\). El operador \(x(n)\) es la derivación
de \(M(1)\) determinada por
\begin{equation}
 x(n)(y\,t^{-m})=n\,\delta_{n,m}\langle x,y\rangle.
 \label{eq:df-derivacion-oscilatoria}
\end{equation}
Ambos actúan como la identidad sobre el segundo factor de
\eqref{eq:df-espacio-estados}. El modo cero actúa sobre la carga mediante
\[
 x(0)(v\otimes\mathbf e^y)
 =\langle x,y\rangle v\otimes\mathbf e^y.
\]

\begin{theorem}[Relaciones de conmutación y no vacuidad]
\label{thm:df-heisenberg}
Para \(m,n\geq1\),
\begin{align}
 [x(n),y(-m)]&=n\delta_{n,m}\langle x,y\rangle I,
 \label{eq:df-CCR}\\
 [x(-n),y(-m)]&=0,
 & [x(n),y(m)]&=0.
\end{align}
Además, \(x(n)\mathbf1=0\) y \(\mathbf1\ne0\).
\end{theorem}

\begin{proof}
Si \(D\) es una derivación y \(M_a\) la multiplicación por \(a\), la
regla de Leibniz da \([D,M_a](v)=D(a)v\). Aplicada a
\eqref{eq:df-derivacion-oscilatoria}, produce \eqref{eq:df-CCR}.
Los multiplicadores conmutan porque el álgebra simétrica es conmutativa;
las derivaciones de coordenadas conmutan en los generadores y, por
Leibniz, en todos los monomios. La anulación del vacío se obtiene al
derivar la unidad. Finalmente, el funcional que toma el coeficiente del
monomio vacío y de la carga cero vale \(1\) sobre \(\mathbf1\), de modo
que el vacío es distinto de cero. La extensión por producto tensorial
conserva todas las identidades.
\end{proof}

Cada monomio oscilatorio tiene un peso entero
\(\operatorname{wt}(a)=\sum_{n,i}n a_{n,i}\), donde
\(a_{n,i}\) son sus ocupaciones. Esta graduación proporciona los límites
finitos de las sumas que aparecen en los coeficientes de campo.

\subsection{Coeficientes de creación y aniquilación en todos los grados}
\label{subsec:df-exponenciales}

% Fuentes: deltas/campos/LatticeCreationExponential.lean,
% LatticeAnnihilationExponential.lean y LatticeAnnihilationEnergy.lean.

Se construyen las dos series formales de operadores
\begin{equation}
 E^-_x(z)=\exp\!\left(\sum_{n\geq1}\frac{x(-n)}n z^n\right)
         =\sum_{d\geq0}C_x(d)z^d,
 \label{eq:df-creacion}
\end{equation}
\begin{equation}
 E^+_x(s)=\exp\!\left(-\sum_{n\geq1}\frac{x(n)}n s^n\right)
         =\sum_{r\geq0}A_x(r)s^r.
 \label{eq:df-aniquilacion}
\end{equation}
Estas exponenciales son operaciones formales sobre los modos ya
construidos. Cada coeficiente se define por una suma finita. En particular,
si \(P_x(z)=\sum_{n\geq1}x(-n)z^n/n\), entonces
\begin{equation}
 C_x(d)=\sum_{q=0}^{d}\frac1{q!}[z^d]P_x(z)^q,
 \label{eq:df-coeficiente-creacion}
\end{equation}
y se obtiene la fórmula análoga para \(A_x(r)\) sustituyendo el potencial
de aniquilación. El orden de composición de los endomorfismos se conserva
en cada potencia.

\begin{lemma}[Estabilización de los coeficientes]
\label{lem:df-estabilizacion}
Para cualquier potencial formal \(P\) con término constante cero,
\([z^d]P^q=0\) cuando \(q>d\). Por tanto, ampliar el límite superior
de \eqref{eq:df-coeficiente-creacion} a cualquier \(N\geq d\) conserva
su valor. Además,
\[
 C_x(0)=A_x(0)=I,
 \qquad A_x(r)1=0\quad(r>0).
\]
Si \(v\) tiene peso a lo sumo \(N\), entonces \(A_x(r)v=0\) para
\(r>N\).
\end{lemma}

\begin{proof}
Cada factor de \(P\) tiene grado al menos uno, por lo que un producto de
\(q\) factores tiene grado al menos \(q\). Este argumento utiliza
únicamente la graduación y se aplica también a coeficientes no
conmutativos. Las identidades en grado cero provienen de la potencia
vacía. Toda composición no vacía de modos de aniquilación anula la
unidad. Finalmente, un término de grado \(r\) en el potencial de
aniquilación reduce el peso oscilatorio en \(r\). Una reducción mayor
que el peso inicial da cero; por linealidad, la conclusión se conserva
en todo el subespacio de peso a lo sumo \(N\).
\end{proof}

Los primeros coeficientes ilustran la normalización:
\begin{align*}
 C_x(1)&=x(-1),&
 C_x(2)&=\tfrac12\bigl(x(-1)^2+x(-2)\bigr),\\
 A_x(1)&=-x(1),&
 A_x(2)&=\tfrac12\bigl(x(1)^2-x(2)\bigr).
\end{align*}
Cada denominador procede de la frecuencia o del orden de la potencia
formal. La forma de Gram que entra en los conmutadores sigue siendo la
forma entera del retículo marcado.

\subsection{Campo de carga y cota de truncamiento de Laurent}
\label{subsec:df-campo-carga}

% Fuente: deltas/campos/LatticeChargedVertexField.lean.

El campo asociado a una carga \(x\in\Lambda_o\) se escribe
\begin{equation}
 Y_x(z)=E^-_x(z)E^+_x(z^{-1})\,e_x z^{x(0)},
 \qquad Y_x(z)=\sum_{k\in\mathbb Z}Y_x[k]z^k,
 \label{eq:df-campo}
\end{equation}
donde \(e_x\mathbf e^y=\varepsilon(x,y)\mathbf e^{x+y}\).
La siguiente fórmula define sus coeficientes sobre una base monomial y
explicita la finitud de cada operación.

\begin{definition}[Coeficiente de campo sobre un estado cargado]
Si \(v\) es un monomio de peso \(N\), se define
\begin{equation}
 Y_x[k](v\otimes\mathbf e^y)
 =\varepsilon(x,y)
 \sum_{j=0}^{N}
 C_x\bigl(k-\langle x,y\rangle+j\bigr)A_x(j)v
       \otimes\mathbf e^{x+y},
 \label{eq:df-coeficiente-campo}
\end{equation}
con la convención \(C_x(d)=0\) para \(d<0\). Se extiende linealmente
a \(V_{\Lambda_o}\).
\end{definition}

\begin{theorem}[Estabilidad, truncamiento y creación de estados de carga]
\label{thm:df-campo-truncamiento}
La fórmula \eqref{eq:df-coeficiente-campo} es independiente de cualquier
cota de peso superior a \(N\). Se cumple
\begin{equation}
 k<\langle x,y\rangle-N
 \quad\Longrightarrow\quad
 Y_x[k](v\otimes\mathbf e^y)=0.
 \label{eq:df-truncamiento-laurent}
\end{equation}
Cada estado admite, por tanto, una cota inferior de Laurent. Sobre el
vacío,
\begin{equation}
 Y_x[k]\mathbf1=
 \begin{cases}
 C_x(k)1\otimes\mathbf e^x,&k\geq0,\\
 0,&k<0.
 \end{cases}
 \label{eq:df-campo-vacio}
\end{equation}
En particular, \(Y_x[0]\mathbf1=1\otimes\mathbf e^x\).
\end{theorem}

\begin{proof}
La ampliación de la cota añade términos \(A_x(j)v\) con \(j>N\), que
son cero por el lema~\ref{lem:df-estabilizacion}. Si se cumple la
desigualdad de \eqref{eq:df-truncamiento-laurent}, todos los índices
\(k-\langle x,y\rangle+j\), con \(j\leq N\), son negativos; cada
sumando se anula. Una combinación lineal finita de monomios conserva
la menor de sus cotas, lo que prueba la afirmación para cualquier
estado. Sobre \(\mathbf1\), sólo subsiste \(j=0\), y
\(\varepsilon(x,0)=1\), \(\langle x,0\rangle=0\). Se obtiene
\eqref{eq:df-campo-vacio}.
\end{proof}

\begin{theorem}[Generación de todo el espacio reticular de estados]
\label{thm:df-generacion-campos}
El subespacio generado a partir del vacío por composiciones finitas de
los operadores \(Y_x[0]\), \(x\in\Lambda_o\), y de todos los modos
de creación \(b_i(-n)\), \(n\geq1\), es \(V_{\Lambda_o}\).
\end{theorem}

\begin{proof}
Por \eqref{eq:df-campo-vacio}, \(Y_x[0]\) aplicado al vacío produce
el estado fundamental de cualquier carga \(x\). Los operadores de
creación multiplican por los generadores del álgebra simétrica;
composiciones finitas de ellos producen todos los monomios sobre ese
estado fundamental. Los monomios oscilatorios por los elementos
\(\mathbf e^x\) constituyen una base algebraica de
\eqref{eq:df-espacio-estados}. Su envolvente lineal es, por tanto, todo
el espacio. El argumento comprende cualquier frecuencia, ocupación
finita y carga reticular.
\end{proof}

\begin{corollary}[Determinación de los entrelazadores por el vacío]
Dos endomorfismos de \(V_{\Lambda_o}\) que conmutan con todos los
generadores del teorema anterior y coinciden sobre el vacío son iguales.
\end{corollary}

\begin{proof}
La conmutación transporta la igualdad en el vacío a cualquier
composición finita de generadores. La generación total y la linealidad
extienden esa igualdad a todo el espacio.
\end{proof}

\subsection{Contracción escalar y ordenación normal de los productos}
\label{subsec:df-contraccion}

% Fuentes: deltas/generacion/LatticeExponentialContraction.lean,
% LatticeContractionFactor.lean, LatticeNormalOrdering.lean;
% deltas/productos/LatticeVacuumContraction.lean.

El producto entero \(p=\langle x,y\rangle\) fija una sucesión escalar
por la recurrencia
\begin{equation}
 s_p(0)=1,\qquad
 (r+1)s_p(r+1)=(r-p)s_p(r).
 \label{eq:df-contraccion-recurrencia}
\end{equation}
La recurrencia determina cada coeficiente a partir del anterior. Su
expresión cerrada es
\(s_p(r)=(-1)^r\binom pr\), con el binomial entero generalizado.
Los primeros términos son
\[
 1,\quad -p,\quad \frac{p(p-1)}2,\quad
 -\frac{p(p-1)(p-2)}6.
\]

\begin{proposition}[Factor formal de contracción]
\label{prop:df-factor-contraccion}
La serie \(S_p(u)=\sum_{r\geq0}s_p(r)u^r\) satisface
\begin{equation}
 (1-u)S_p'(u)=-pS_p(u),\qquad S_p(0)=1.
 \label{eq:df-ecuacion-factor}
\end{equation}
Es la única solución formal de esas condiciones. Además,
\[
 S_{p+q}=S_pS_q,qquad S_0=1,qquad S_1=1-u,qquad
 S_pS_{-p}=1.
\]
Para \(p\geq0\), \(S_p(u)=(1-u)^p\), y sus coeficientes de grado
superior a \(p\) se anulan.
\end{proposition}

\begin{proof}
Comparar el coeficiente de \(u^r\) en
\eqref{eq:df-ecuacion-factor} reproduce exactamente
\eqref{eq:df-contraccion-recurrencia}. Como \(r+1\ne0\), los
coeficientes quedan determinados inductivamente. La regla de Leibniz
muestra que \(S_pS_q\) satisface la misma ecuación que \(S_{p+q}\)
y tiene el mismo término constante; la unicidad da la identidad
multiplicativa. Los casos \(0,1\) se verifican directamente. La
identidad inversa resulta del caso \(q=-p\), y las potencias naturales
se obtienen por inducción. Si \(p\) es natural, el paso \(r=p\) de
la recurrencia produce cero y los pasos posteriores conservan ese cero.
\end{proof}

\begin{theorem}[Ordenación normal en todos los grados]
\label{thm:df-ordenacion}
Las series construidas a partir de los modos satisfacen
\begin{equation}
 E^+_x(s)E^-_y(z)
 =S_{\langle x,y\rangle}(sz),E^-_y(z)E^+_x(s).
 \label{eq:df-ordenacion-normal}
\end{equation}
En coeficientes, para \(r,d\geq0\),
\begin{equation}
 A_x(r)C_y(d)
 =\sum_{i=0}^{\min(r,d)}
 s_{\langle x,y\rangle}(i),
 C_y(d-i)A_x(r-i).
 \label{eq:df-ordenacion-coeficientes}
\end{equation}
\end{theorem}

\begin{proof}
Sean
\(Q_x(s)=-\sum_{n\geq1}x(n)s^n/n\) y
\(P_y(z)=\sum_{m\geq1}y(-m)z^m/m\).
Por \eqref{eq:df-CCR},
\[
 [Q_x(s),P_y(z)]
 =-\langle x,y\rangle\sum_{n\geq1}\frac{(sz)^n}{n}\,I.
\]
El conmutador es central. La identidad exponencial para dos potenciales
con conmutador central se demuestra aquí coeficiente a coeficiente:
conmutar una potencia de \(Q_x\) a través de una potencia de \(P_y\)
añade los términos escalares de ese conmutador; las sumas a cada
bigrado son finitas por el lema~\ref{lem:df-estabilizacion}. El factor
escalar resultante tiene término constante uno y satisface
\eqref{eq:df-ecuacion-factor}, con
\(p=\langle x,y\rangle\); por la unicidad de la
proposición~\ref{prop:df-factor-contraccion}, es \(S_p(sz)\).
Esto demuestra \eqref{eq:df-ordenacion-normal}. La extracción del
coeficiente \(s^rz^d\) da
\eqref{eq:df-ordenacion-coeficientes}; el factor de contracción consume
el mismo grado \(i\) en ambas variables, de donde procede el límite
\(i\leq\min(r,d)\).
\end{proof}

\begin{theorem}[Contracción exacta de los estados creados desde el vacío]
\label{thm:df-contraccion-vacio}
Para toda pareja de cargas \(x,y\) y todos los grados \(r,d\),
\begin{equation}
 A_x(r)C_y(d)1=
 \begin{cases}
 s_{\langle x,y\rangle}(r),C_y(d-r)1,&r\leq d,\\
 0,&r>d.
 \end{cases}
 \label{eq:df-contraccion-vacio}
\end{equation}
En la diagonal \(r=d\), el resultado es
\(s_{\langle x,y\rangle}(r)1\). Si
\(\langle x,y\rangle=p\geq0\), también se anula para \(r>p\).
\end{theorem}

\begin{proof}
Evaluamos \eqref{eq:df-ordenacion-coeficientes} sobre \(1\).
El lema~\ref{lem:df-estabilizacion} anula todos los sumandos con
\(r-i>0\). Subsiste únicamente \(i=r\), siempre que \(r\leq d\).
Para \(r=d\), \(C_y(0)=I\). La última anulación sigue de la cota
polinómica de \(S_p\).
\end{proof}

Por ejemplo, \(A_x(1)C_y(1)1=-\langle x,y\rangle1\).
Para \(p=-1\), la recurrencia da \(s_{-1}(r)=1\) en todos los grados;
para \(p=2\), da \((1,-2,1,0,\ldots)\). Estas dos evaluaciones
exhiben, respectivamente, una contracción de orden ilimitado y una
cancelación polinómica finita, ambas determinadas por la misma forma
entera del retículo.

\subsection{Conmutador mixto en todos los modos enteros}
\label{subsec:df-conmutador-mixto}

% Fuentes: deltas/integracion/LatticePositiveMixedFields.lean,
% LatticeNegativeMixedFields.lean y LatticeAllMixedFields.lean.

Para \(x\in\Lambda_o\), escribimos \(H_x(m)=x(m)\), con
\(m\in\mathbb Z\), incluyendo los modos negativos de creación,
el modo cero de carga y los positivos de aniquilación.

\begin{theorem}[Conmutador mixto integral]
\label{thm:df-conmutador-mixto}
Para todas las cargas \(x,y\) y todos los enteros \(m,k\),
\begin{equation}
 [H_x(m),Y_y[k]]
 =\langle x,y\rangle,Y_y[k-m].
 \label{eq:df-conmutador-mixto}
\end{equation}
\end{theorem}

\begin{proof}
Tratamos separadamente los tres signos de \(m\), conservando en cada
caso los coeficientes de campo de \eqref{eq:df-coeficiente-campo}.

Si \(m>0\), la relación de Heisenberg da
\[
 [x(m),C_y(d)]=
 \begin{cases}
 \langle x,y\rangle C_y(d-m),&d\geq m,\\
 0,&d<m.
 \end{cases}
\]
Los modos positivos conmutan con los coeficientes de aniquilación y
con el desplazamiento de carga. Al conmutar término a término en
\eqref{eq:df-coeficiente-campo}, el grado de creación disminuye en
\(m\), que equivale a reemplazar \(k\) por \(k-m\). El corte de peso
permanece válido porque la aniquilación disminuye el peso.

Si \(m=0\), el campo desplaza la carga \(z\) a \(y+z\). La diferencia
de los autovalores del modo cero es
\(\langle x,y+z\rangle-\langle x,z\rangle=\langle x,y\rangle\).
Los operadores oscilatorios preservan la carga, de modo que se obtiene
\eqref{eq:df-conmutador-mixto} con \(k-m=k\).

Finalmente, sea \(m=-a\), con \(a>0\). Los modos de creación conmutan
entre sí, mientras que la relación derivada de
\eqref{eq:df-CCR} y \eqref{eq:df-aniquilacion} es
\[
 [x(-a),A_y(j)]=
 \begin{cases}
 \langle x,y\rangle A_y(j-a),&j\geq a,\\
 0,&j<a.
 \end{cases}
\]
Si el estado inicial tiene peso a lo sumo \(N\), después de crear el
modo \(-a\) tiene peso a lo sumo \(N+a\). Usamos este corte común en
los dos términos del conmutador. Los sumandos con \(j<a\) se anulan;
reindexamos los restantes mediante \(j=r+a\), \(0\leq r\leq N\).
El grado de creación se transforma en
\(k-\langle y,z\rangle+r+a\), exactamente el coeficiente de campo
con índice \(k+a=k-m\). La identidad queda probada sobre los estados
monomiales cargados y se extiende por linealidad a todo el espacio.
\end{proof}

El desplazamiento del índice de campo distingue correctamente la
orientación de los modos: un modo positivo resta su frecuencia, un modo
negativo la suma y el modo cero conserva el índice. Los tres casos son
realizaciones de una única identidad en \(\mathbb Z\).

\subsection{Localidad mixta por cancelación coeficiente a coeficiente}
\label{subsec:df-localidad-mixta}

% Fuentes: deltas/integracion/LatticeMixedLocalityCriterion.lean,
% LatticeMixedLocality.lean y SelectedMixedFieldLocality.lean.

Consideremos el campo de Heisenberg
\(H_x(z)=\sum_{m\in\mathbb Z}H_x(m)z^{-m-1}\).
Para una serie bivariada de coeficientes \(F(a,b)\), la multiplicación
por \(z-w\) actúa como
\begin{equation}
 (\mathcal C F)(a,b)=F(a-1,b)-F(a,b-1).
 \label{eq:df-operador-cruce}
\end{equation}
Las composiciones ordenadas de los dos campos tienen coeficientes
\[
 F(a,b)=H_x(-a-1)Y_y[b],\qquad
 G(a,b)=Y_y[b]H_x(-a-1).
\]

\begin{theorem}[Localidad mixta de orden uno]
\label{thm:df-localidad-mixta}
Los campos construidos satisfacen
\begin{equation}
 (z-w)H_x(z)Y_y(w)=(z-w)Y_y(w)H_x(z).
 \label{eq:df-localidad-mixta}
\end{equation}
La igualdad es una identidad de coeficientes endomorfos en todo
\(V_{\Lambda_o}\).
\end{theorem}

\begin{proof}
Por \eqref{eq:df-operador-cruce}, el coeficiente \((a,b)\) de la
diferencia es
\[
 [H_x(-a),Y_y[b]]-[H_x(-a-1),Y_y[b-1]].
\]
El teorema~\ref{thm:df-conmutador-mixto} convierte ambos términos en
\(\langle x,y\rangle Y_y[a+b]\). Su resta es cero para todos los
enteros \(a,b\). La identidad se cumple en los endomorfismos, por lo
que vale simultáneamente para todos los estados.
\end{proof}

\subsection{Localidad de los campos cargados y conservación del origen}
\label{subsec:df-localidad-cargada}

% Fuentes: deltas/localidad/LatticeFactorConvolution.lean,
% LatticeChargedLocality.lean, LatticeChargedLocalityFull.lean,
% SelectedFieldLocality.lean y deltas/integracion/SelectedMixedFieldLocality.lean.

La ordenación normal permite demostrar también la localidad de dos
campos de carga. Sea \(p=\langle x,y\rangle\) y
\(N_p=\max(0,-p)\). El factor de contracción de los dos órdenes se
expande en las regiones formales \(w/z\) y \(z/w\); su paridad relativa
está fijada por \eqref{eq:df-conmutador-cociclo}.

\begin{theorem}[Localidad uniforme de los campos cargados]
\label{thm:df-localidad-cargada}
Para cualesquiera \(x,y\in\Lambda_o\),
\begin{equation}
 (z-w)^{N_p}Y_x(z)Y_y(w)
 =(z-w)^{N_p}Y_y(w)Y_x(z),
 \qquad N_p=\max(0,-\langle x,y\rangle).
 \label{eq:df-localidad-cargada}
\end{equation}
El exponente depende de las cargas y es válido para todos los estados.
\end{theorem}

\begin{proof}
Sobre un estado de peso acotado, los productos de los coeficientes de
campo se expresan mediante \eqref{eq:df-ordenacion-coeficientes}, los
desplazamientos de carga y el cociclo. Las sumas por coeficiente son
finitas: la aniquilación tiene corte de peso y las dos variables del
producto ordenado tienen una cota inferior. La identidad cocíclica
reúne los desplazamientos en el mismo elemento de carga \(x+y\), y
la identidad de conmutador aporta el factor \((-1)^p\) entre los
dos órdenes.

Los factores escalares resultantes son las dos expansiones de
\((z-w)^p\). Si \(p\geq0\), ambas son el mismo polinomio y coinciden
sin multiplicación adicional. Si \(p<0\), multiplicar por
\((z-w)^{-p}\) transforma ambas expansiones en el factor uno.
Formalmente, cada aplicación de \(\mathcal C\) eleva en uno el
exponente de ambas expansiones; después de \(-p\) aplicaciones se
alcanza el exponente cero. Esta cancelación es independiente del corte
de peso elegido y de la carga del estado sobre el que actúan los
campos. Se obtiene así \eqref{eq:df-localidad-cargada} sobre todo
estado monomial cargado, y la linealidad la extiende a
\(V_{\Lambda_o}\).
\end{proof}

\begin{theorem}[Composición conservativa de incidencia, generación y localidad]
\label{thm:df-composicion-campos}
Fijado el origen \(o\) producido por la incidencia del registro
seleccionado, el mismo espacio \(V_{\Lambda_o}\) satisface conjuntamente:
la generación desde el vacío del
teorema~\ref{thm:df-generacion-campos}, los productos de todos los grados
del teorema~\ref{thm:df-contraccion-vacio}, la localidad cargada de
\eqref{eq:df-localidad-cargada} y la localidad mixta de
\eqref{eq:df-localidad-mixta}. Su forma entera, su cociclo, su vacío
y sus operadores de carga son los definidos a partir de
\(\Lambda_o\) en esta sección.
\end{theorem}

\begin{proof}
Todos los enunciados anteriores están cuantificados sobre el mismo
retículo y la misma forma entera. Al especializarlos al origen
\(o\) ya seleccionado, sus objetos coinciden por sus definiciones:
el cociclo usa la base de \(\Lambda_o\), los osciladores usan su
Gram, los campos utilizan esos osciladores y ese cociclo, y los
productos utilizan esos mismos coeficientes de campo. La conjunción
de los teoremas preserva, por tanto, la bandera de incidencia de
\eqref{eq:df-soportes-K} y todas las dependencias que determinaron
el origen. Ninguna especialización introduce una nueva elección
reticular.
\end{proof}

La concatenación demostrativa distingue con precisión sus operaciones:
los bloques regionales determinan prefijos, cilindros y firmas; la
selección dodecafásica y la incidencia producen el origen reticular;
el retículo determina la forma entera y el cociclo; éstos fijan los
modos, los campos y sus productos. La estructura de campos resultante
queda así provista de generación total y de identidades de localidad
en el espacio algebraico completo. Los capítulos dedicados a la
reconstrucción de la estructura de operadores de vértice, al sector
torcido y a la realización excepcional utilizan estos resultados con
sus respectivas hipótesis adicionales, manteniendo esta misma
procedencia generativa.
